% =============================================================
%  Overhead lines — what sets R', L', G', C' on an OHL
%  The general representation is Section 3; the worked numbers for
%  the 500 kV reference line are in 06-reference-line.tex.
% =============================================================
\subsection{Overhead Lines}
\label{sec:ohl}

\subsubsection{Geometry and the line parameters}
\label{sec:ohl-geometry}

On an overhead line the conductors are bare and the insulation is air, so
the inductance and the capacitance follow from the geometry of the tower and
of the conductor bundle alone:
\begin{equation}
  L' = \frac{\mu_0}{2 \cdot \pi} \cdot \ln\frac{D_m}{r_b}
  \qquad
  C' = \frac{2 \cdot \pi \cdot \varepsilon_0}{\ln\left(D_m / r_{b\_c}\right)}
  \label{eq:ohl-lc}
\end{equation}
with $D_m$ the geometric mean distance between the phases and $r_b$,
$r_{b\_c}$ the equivalent radii of the bundle for the inductance and the
capacitance respectively. The two differ because the inductance is referred
to the geometric mean radius of the sub-conductors and the capacitance to
their physical radius.

Both parameters depend on the geometry only through a logarithm, and on the
same ratio $D_m / r_b$. That makes the geometry a weak lever: the phase
spacing and the bundle radius have to be changed by a factor to move $L'$ or
$C'$ by a fraction. Length, by contrast, enters linearly, so on an overhead
line the series reactance of~\eqref{eq:ac-x} is set by the route far more
than by the tower.

The flat (horizontal) configuration used throughout this report is drawn to
scale in Figure~\ref{fig:ohl-tower}: three phases on one crossarm at a
spacing $D$, so that the outer phases see $2 \cdot D$, each phase a bundle
of three sub-conductors. Table~\ref{tab:ohl-geometry} gives its dimensions.
It gives the lowest towers and the widest right-of-way of the common
arrangements.

\begin{figure}[!ht]
  \centering
  \begin{tikzpicture}[x=1.15mm, y=1.15mm,
      cond/.style={fill=atFg, draw=atFg},
      steel/.style={atBorderStrong},
      dim/.style={atFgMuted, {Stealth[length=1.6mm]}-{Stealth[length=1.6mm]},
                  line width=0.5pt},
      lbl/.style={font=\footnotesize, atFg},
      note/.style={font=\scriptsize, atFgMuted}]
    % crossarm and tower body, schematic
    \draw[steel, line width=1.1pt] (-40,0) -- (40,0);
    \draw[steel, line width=1.1pt] (-40,0) -- (-38,-3) -- (38,-3) -- (40,0);
    \draw[steel, line width=0.9pt] (-6,0) -- (-7.8,-18);
    \draw[steel, line width=0.9pt] ( 6,0) -- ( 7.8,-18);
    \foreach \yy in {-6,-12} {
      \draw[steel, line width=0.4pt] ({-6-abs(\yy)/10},\yy) -- ({6+abs(\yy)/10},\yy);
    }
    % insulator strings and bundles
    \foreach \xx in {-35,0,35} {
      \draw[atFgMuted, line width=0.8pt] (\xx,0) -- (\xx,-12);
      \foreach \yy in {-2,-4,...,-10} {\fill[atFgMuted] (\xx,\yy) circle (0.7);}
      \foreach \px/\py in {0/-0.87,-0.75/0.43,0.75/0.43} {
        \fill[cond] ($(\xx,-13)+(\px,\py)$) circle (0.55);
      }
    }
    \node[lbl] at (-35,-17) {A};
    \node[lbl] at (  0,-17) {B};
    \node[lbl] at ( 35,-17) {C};
    % extension lines
    \draw[atFgMuted, line width=0.4pt, densely dotted] (-35,-15.5) -- (-35,-30);
    \draw[atFgMuted, line width=0.4pt, densely dotted] (  0,-15.5) -- (  0,-24);
    \draw[atFgMuted, line width=0.4pt, densely dotted] ( 35,-15.5) -- ( 35,-30);
    % dimensions
    \draw[dim] (-35,-22) -- (0,-22) node[midway,below=1pt,lbl] {$D = 35$\,ft};
    \draw[dim] (0,-22) -- (35,-22) node[midway,below=1pt,lbl] {$D = 35$\,ft};
    \draw[dim] (-35,-28) -- (35,-28) node[midway,below=1pt,lbl] {$2 \cdot D = 70$\,ft};
    % detail marker on phase C
    \draw[atAccentFg, line width=0.8pt] (35,-13) circle (2.4);
    \node[note] at (0,4) {to scale, 1\,ft = 1.15\,mm; earth return neglected};
  \end{tikzpicture}
  \caption{Tower head of the reference line, flat configuration; the marked
           phase is drawn in detail in Figure~\ref{fig:ohl-bundle-detail}.
           Tower body and insulator lengths are schematic}
  \label{fig:ohl-tower}
\end{figure}

\begin{table}[!ht]
  \centering
  \begin{tabular}{@{}l l r@{}}
    \toprule
    \textbf{Parameter} & & \multicolumn{1}{c}{\textbf{Value}} \\
    \midrule
    $D$ phase spacing & ft & 35 \\
    $2 \cdot D$ outer phases & ft & 70 \\
    Sub-conductors per phase & & 3 \\
    $s$ bundle spacing & in & 18 \\
    $d$ conductor diameter & in & 1.762 \\
    $GMR$ of a sub-conductor & ft & 0.0588 \\
    \bottomrule
  \end{tabular}
  \caption{Tower-head geometry of the reference line, Bluebird ACSR
           2156\,kcmil 84/19, single circuit~[10]}
  \label{tab:ohl-geometry}
\end{table}

Bundling is the one geometric change that is worth making for its electrical
effect. Dividing the phase current between sub-conductors raises the
equivalent radius $r_b$ well above that of a single conductor, and with it
reduces $L'$; it also reduces the surface gradient and so the corona loss.
Table~\ref{tab:ohl-bundle} shows the effect for one configuration: going
from one sub-conductor to three removes a third of the reactance, while
doubling $r_b$ again --- or halving $D_m$ --- removes only a further
15\,\%.

\begin{table}[!ht]
  \centering
  \begin{tabular}{@{}c r r r@{}}
    \toprule
    \textbf{Sub-conductors per phase} & \multicolumn{1}{c}{\textbf{$r_b$}}
      & \multicolumn{1}{c}{\textbf{$X'$}}
      & \multicolumn{1}{c}{\textbf{relative}} \\
    & \multicolumn{1}{c}{m} & \multicolumn{1}{c}{$\Omega$/km}
      & \multicolumn{1}{c}{\%} \\
    \midrule
    1 & 0.018 & 0.499 & 148 \\
    2 & 0.091 & 0.377 & 112 \\
    3 & 0.155 & 0.336 & 100 \\
    4 & 0.222 & 0.310 &  92 \\
    \bottomrule
  \end{tabular}
  \caption{Effect of bundling on the series reactance, 35\,ft flat
           configuration with 18\,in bundle spacing at 60\,Hz}
  \label{tab:ohl-bundle}
\end{table}

The shunt branch of an overhead line is effectively a pure capacitance. The
conductance $G'$ collects the leakage across the insulator strings and the
corona loss, and in normal weather it is two orders of magnitude below the
susceptance $\omega \cdot C'$, so it is set to zero in load-flow work. It is
not negligible for loss studies in foul weather, where corona dominates it.

Because the conductors are far apart and surrounded by air, $C'$ stays
small. The wave impedance $\underline{Z}_w = \sqrt{\underline{Z}' /
\underline{Y}'}$ is correspondingly high --- a few hundred ohms --- and the
charging current per unit length is low enough that an overhead line can run
for hundreds of kilometres before it has to be compensated.

\subsubsection{Bundle and phase equivalents}
\label{sec:ohl-bundle}

The two radii of~\eqref{eq:ohl-lc} and the mean phase distance follow from
the tower head and the bundle. For a three-conductor bundle of spacing $s$
made of sub-conductors of diameter $d$ and geometric mean radius $GMR$,
Figure~\ref{fig:ohl-bundle-detail}, and a flat configuration of phase
spacing $D$ in which the outer phases see $2 \cdot D$,
\begin{equation}
  GMR_b = \sqrt[3]{GMR \cdot s^{2}},
  \qquad
  r_b = \sqrt[3]{\frac{d}{2} \cdot s^{2}},
  \qquad
  GMD = \sqrt[3]{D \cdot D \cdot 2 \cdot D} = D \cdot \sqrt[3]{2}
  \label{eq:ohl-equivalents}
\end{equation}
\begin{figure}[!ht]
  \centering
  \begin{tikzpicture}[x=1.65mm, y=1.65mm,
      dim/.style={atFgMuted, {Stealth[length=1.4mm]}-{Stealth[length=1.4mm]},
                  line width=0.5pt},
      lbl/.style={font=\footnotesize, atFg},
      note/.style={font=\scriptsize, atFgMuted}]
    % equilateral triangle, s = 18 in, apex down; circumradius 10.39 in
    \coordinate (P1) at (0,-10.39);
    \coordinate (P2) at (-9,5.196);
    \coordinate (P3) at ( 9,5.196);
    \draw[atBorderStrong, line width=0.5pt, densely dashed] (P1) -- (P2) -- (P3) -- cycle;
    \draw[atBorderStrong, line width=0.4pt, densely dotted] (0,0) circle (10.39);
    \foreach \P in {P1,P2,P3} {
      \fill[atFg] (\P) circle (0.881);
    }
    \fill[atFgMuted] (0,0) circle (0.35);
    \draw[dim] ($(P2)+(0,4)$) -- ($(P3)+(0,4)$) node[midway,above,lbl] {$s = 18$\,in};
    \draw[atFgMuted, line width=0.4pt, densely dotted] (P2) -- ($(P2)+(0,4)$);
    \draw[atFgMuted, line width=0.4pt, densely dotted] (P3) -- ($(P3)+(0,4)$);
    \draw[dim] ($(P1)+(-0.881,-2.8)$) -- ($(P1)+(0.881,-2.8)$);
    \node[lbl,anchor=west] at ($(P1)+(1.8,-2.8)$) {$d = 1.762$\,in};
    \draw[atFgMuted, line width=0.4pt, densely dotted] ($(P1)+(-0.881,0)$) -- ($(P1)+(-0.881,-2.8)$);
    \draw[atFgMuted, line width=0.4pt, densely dotted] ($(P1)+( 0.881,0)$) -- ($(P1)+( 0.881,-2.8)$);
    \draw[dim] (0,0) -- (P3) node[midway,above left,lbl,xshift=-1mm] {$s/\sqrt{3}$};
    \node[note] at (0,-16) {to scale, 1\,in = 1.9\,mm};
  \end{tikzpicture}
  \caption{Three-conductor bundle of the reference line, apex down, with the
           circle of radius $s/\sqrt{3}$ on which the sub-conductors sit}
  \label{fig:ohl-bundle-detail}
\end{figure}

with $GMR_b$ the radius entering $L'$ and $r_b$ the one entering $C'$. For a
two-conductor bundle the same construction gives $GMR_b = \sqrt{GMR \cdot
s}$ and $r_b = \sqrt{d/2 \cdot s}$. The $GMD$ above assumes a transposed
line.

Table~\ref{tab:ohl-equivalents} evaluates them for the reference tower head
of Section~\ref{sec:reference-line} --- 35\,ft flat, 18\,in bundle spacing,
Bluebird ACSR 2156\,kcmil with $d = 1.762$\,in and $GMR = 0.0588$\,ft ---
with three and with two sub-conductors per phase.

\begin{table}[!ht]
  \centering
  \begin{tabular}{@{}l l r r@{}}
    \toprule
    \textbf{Quantity} & & \multicolumn{1}{c}{\textbf{$3\times$}}
      & \multicolumn{1}{c}{\textbf{$2\times$}} \\
    \midrule
    $GMR_b$ & in & 6.11 & 3.56 \\
    $r_b$   & in & 6.58 & 3.98 \\
    $GMD$   & ft & 44.1 & 44.1 \\
    \midrule
    $L'$ & mH/km & 0.892 & 1.000 \\
    $C'$ & nF/km & 12.68 & 11.38 \\
    $Z_w = \sqrt{L'/C'}$ & $\Omega$ & 265 & 296 \\
    \bottomrule
  \end{tabular}
  \caption{Bundle equivalents and the resulting parameters, three against
           two sub-conductors per phase}
  \label{tab:ohl-equivalents}
\end{table}

The third sub-conductor lowers $Z_w$ by about 10\,\% and so raises the
natural power of Section~\ref{sec:natural-power} by about 11\,\%. A flat
configuration gives the lowest towers and the widest right-of-way; the outer
phases see twice the phase spacing, which is what the cube root
in~\eqref{eq:ohl-equivalents} averages.

\subsubsection{Conductor codewords}
\label{sec:ohl-codewords}

A conductor size does not define a conductor. ACSR of 795\,kcmil exists in
four standard strandings --- Cuckoo 24/7, Drake 26/7, Mallard 30/19 and
Condor 54/7 --- each with a different steel content and therefore different
strength, weight and sag. The codeword fixes area, stranding and core in one
word.

The names began as Alcoa product nomenclature in the 1920s and 1930s and
were later adopted into the ASTM standards~[13]. The family identifies the
construction: ACSR and ACSS take bird names (Drake, Falcon, Rail, Cardinal),
AAC flowers (Rose, Iris, Tulip, Peony)~[14] and AAAC US city names (Akron,
Amherst, Anaheim)~[15]. The name carries no size information --- Bluebird is
2156\,kcmil, Condor only 795.

Other systems use the same idea with different animals. BS 215 and Indian
practice use mammals: Dog, Wolf, Lynx, Panther, Zebra, Moose~[16]. A 400\,kV
Indian line specified as ``twin Moose'' is the equivalent of a North
American ``bundled Drake''. Metric designs in continental Europe usually
drop codewords altogether and quote the aluminium and steel areas directly,
for example 490/65\,mm$^2$.

The practical consequence is that a conductor has to be quoted with its
standard as well as its codeword, since the same name can appear in more
than one national system.

\FloatBarrier
